\documentclass[10pt]{article}
\usepackage[T1]{fontenc}
\usepackage[utf8]{inputenc}
\usepackage[french]{babel}
\usepackage{amsmath,amssymb,mathrsfs,mathpazo}
\usepackage{geometry}
\geometry{a4paper,margin=2cm}
\usepackage{graphicx,array,tabularx,tabularray,colortbl,multirow,multicol,booktabs,enumitem,pifont,ulem,fancybox,diagbox,cancel}
\usepackage{pstricks,pstricks-add,pst-plot,pst-tree,pst-node,pst-text,pst-func,pst-math,pst-bspline,pst-3dplot,pst-eucl,pst-solides3d}
\let\pscalculate\undefined
\usepackage{pst-bezier}
\usepackage{tikz,pgfplots,tkz-tab}
\usetikzlibrary{arrows,arrows.meta,patterns,calc,positioning,trees,decorations.markings,shapes.geometric,shapes.misc,shapes.arrows,intersections,decorations.pathmorphing,decorations.pathreplacing,backgrounds,shadows,babel}
\pgfplotsset{compat=1.18}
\usepgfplotslibrary{fillbetween}
\usepackage[np]{numprint}
\usepackage{listings,makecell,pas-tableur,twoopt}
\let\barre\undefined
\usepackage[pstricks]{tablvar}
\definecolor{darkgreen}{RGB}{0,100,0}
\providecommand{\red}{\color{red}}
\providecommand{\blue}{\color{blue}}
\providecommand{\green}{\color{green}}
\providecommand{\black}{\color{black}}
\providecommand{\white}{\color{white}}
\providecommand{\gray}{\color{gray}}
\providecommand{\cyan}{\color{cyan}}
\providecommand{\magenta}{\color{magenta}}
\providecommand{\yellow}{\color{yellow}}
\providecommand{\violet}{\color{violet}}
\providecommand{\orange}{\color{orange}}
\providecommand{\euro}{\texteuro}
\providecommand{\EUR}{\texteuro}
\providecommand{\ieme}{\textsuperscript{e}}
\providecommand{\degres}{\ensuremath{^\circ}}
\providecommand{\degree}{\ensuremath{^\circ}}
\providecommand{\vv}[1]{\overrightarrow{#1}}
\providecommand{\Coord}[1]{#1}
\providecommand{\decosix}{}
\providecommand{\starredbullet}{\ensuremath{\bullet}}
\pagestyle{empty}
\setlength{\parindent}{0pt}
\usepackage{adjustbox}\DeclareUnicodeCharacter{2236}{:}\DeclareUnicodeCharacter{2190}{\ensuremath{\leftarrow}}
\providecommand{\R}{\mathbb{R}}
\providecommand{\N}{\mathbb{N}}
\providecommand{\Z}{\mathbb{Z}}
\providecommand{\Q}{\mathbb{Q}}
\providecommand{\D}{\mathbb{D}}
\providecommand{\C}{\mathbb{C}}
\providecommand{\e}{{\mathrm{e}}}
\providecommand{\vect}[1]{\overrightarrow{#1}}
\providecommand{\vectt}[1]{\overrightarrow{\mathrm{#1}}}
\providecommand{\vvt}[1]{\overrightarrow{\mathrm{#1}}}
\providecommand{\np}[2][]{{#2~#1}}
\providecommand{\vv}[1]{{\overrightarrow{#1}}}
\renewcommand{\d}{{\mathrm{d}}}
\providecommand{\degres}{{^\circ}}
\providecommand{\degre}{{^\circ}}
\providecommand{\blue}{\color{blue}}
\providecommand{\red}{\color{red}}
\providecommand{\enskip}{\quad}
\providecommand{\llbracket}{\lbrack\!\lbrack}
\providecommand{\rrbracket}{\rbrack\!\rbrack}
\providecommand{\euro}{{\text{€}}}
\providecommand{\no}{{\text{n°}}}
\providecommand{\interval}[2]{{[#1;#2]}}
\providecommand{\intervFO}[2]{{[#1;#2[}}
\providecommand{\dotfill}{\ldots}

\providecommand{\ds}{\displaystyle}
\providecommand{\textding}[1]{#1}
\providecommand{\Oij}{\(\mathrm{O};\vec{\imath},\vec{\jmath}\)}
\providecommand{\Oijk}{\(\mathrm{O};\vec{\imath},\vec{\jmath},\vec{k}\)}
\providecommand{\Ouv}{\(\mathrm{O};\vec{u},\vec{v}\)}
\providecommand{\intervFF}[2]{\ensuremath{[#1;#2]}}
\renewcommand{\euro}{\ifmmode\text{\texteuro}\else\texteuro\fi}
\DeclareUnicodeCharacter{00B0}{\ifmmode^{\circ}\else\textdegree\fi}
\DeclareUnicodeCharacter{20AC}{\ifmmode\text{\texteuro}\else\texteuro\fi}
% Énoncés : Éducation nationale. Transcription et corrigés : APMEP, Denis Vergès et contributeurs. Découpage et classement : CoopMaths, Sébastien Lozano, Éric Elter et contributeurs. Adaptation de la présentation : Maths Explorer.
% CC BY-SA : https://coopmaths.fr/www/about/
% https://coopmaths.fr/alea/static/e3c/2020/tex/e3c_2020_00_sujet53_2_cor.tex
\begin{document}
\noindent\textit{APMEP / CoopMaths. CC BY-SA. Présentation adaptée par Maths Explorer.}\par\medskip
\subsection*{1.}

On a :
\[
p_T(S) = \dfrac{p(T \cap S)}{p(T)} = \dfrac{0{,}02}{0{,}08} = \dfrac{2}{8} = \dfrac{1}{4} = 0{,}25.
\]

\subsection*{2.}

\begin{center}
\begin{tikzpicture}[thick, scale=1.5]
\node (P_-1_0) at (-2,-1.5) {$\phantom{A}$};
\node (P_0_0) at (0,-0.5) {$T$};
\draw (P_-1_0) -- (P_0_0) node[midway, above] {$0{,}08$};
\node (P_1_0) at (2,-0) {$S$};
\draw (P_0_0) -- (P_1_0) node[midway, above] {$0{,}25$};
\node (P_1_1) at (2,-1) {$\overline{S}$};
\draw (P_0_0) -- (P_1_1) node[midway, below] {$0{,}75$};
\node (P_0_2) at (0,-2.5) {$\overline{T}$};
\draw (P_-1_0) -- (P_0_2) node[midway, below] {$0{,}92$};
\node (P_1_2) at (2,-2) {$S$};
\draw (P_0_2) -- (P_1_2) node[midway, above] {$0{,}05$};
\node (P_1_3) at (2,-3) {$\overline{S}$};
\draw (P_0_2) -- (P_1_3) node[midway, below] {$0{,}95$};
\end{tikzpicture}
\end{center}

\subsection*{3.}

On a :
\begin{itemize}
    \item \(p(\overline{S} \cap T) = p(T \cap \overline{S}) = p(T) \times p_T(\overline{S}) = 0{,}08 \times 0{,}75 = 0{,}06\).
    \item \(p(\overline{S} \cap \overline{T}) = p(\overline{T} \cap \overline{S}) = p(\overline{T}) \times p_{\overline{T}}(\overline{S}) = 0{,}92 \times 0{,}95 = 0{,}874\).
\end{itemize}

D'après la loi des probabilités totales :
\begin{align*}
p(\overline{S}) &= p(\overline{S} \cap T) + p(\overline{S} \cap \overline{T}) \\
&= 0{,}06 + 0{,}874 \\
&= 0{,}934.
\end{align*}

\subsection*{4.}

\paragraph{a.}
\[
\begin{array}{|c|c|c|c|}
\hline
x_i & 0 & 9 & 14 \\
\hline
P(X = x_i) & 0{,}066 & 0{,}06 & 0{,}874 \\
\hline
\end{array}
\]

\paragraph{b.} Le prix de vente moyen d'un jeu fabriqué par cette entreprise est égal à l'espérance mathématique de la variable aléatoire \(X\), soit :
\[
E(X) = 0 \times 0{,}066 + 9 \times 0{,}06 + 14 \times 0{,}874 = 12{,}776,
\]
soit \( 12{,}78 \) € au centime près.

\end{document}
